\chapter*{ОТВЕТЫ НА ВОПРОСЫ}
\addcontentsline{toc}{chapter}{ОТВЕТЫ НА ВОПРОСЫ}

\section*{1. Почему метод Казиски не всегда точен и как повысить его надежность?}

Метод Казиски может давать неточные результаты по нескольким причинам:

\begin{itemize}
    \item повторяющиеся n-граммы могут возникать случайно и не отражать реальную структуру ключа;
    \item текст может быть слишком коротким, чтобы набралось достаточное количество повторений;
    \item если ключ слишком длинный или почти длины текста, метод перестаёт работать;
    \item алфавит и структура языка влияют на частоту естественных повторов.
\end{itemize}

\textbf{Способы повышения надёжности:}

\begin{itemize}
    \item комбинирование метода Казиски с индексом совпадений;
    \item использование более длинных n-грамм (4--5 символов);
    \item фильтрация случайных совпадений по статистической значимости;
    \item проверка нескольких кандидатов длины ключа, а не только максимального НОД.
\end{itemize}

\section*{2. Как длина ключа влияет на сложность взлома шифра Виженера?}

Длина ключа напрямую определяет криптостойкость:

\begin{itemize}
    \item \textbf{Короткий ключ (3--5 символов):} шифротекст легко разобрать на группы с одинаковым сдвигом, и методы Казиски/частотного анализа работают очень точно;
    \item \textbf{Ключ средней длины (8--12 символов):} анализ всё ещё возможен, но требуется больше статистики;
    \item \textbf{Длинный ключ (15+ символов):} число символов в каждой группе становится малым, что ухудшает частотный анализ;
    \item \textbf{Ключ длины текста:} шифр становится одноразовым блокнотом и практически не поддаётся классическому криптоанализу.
\end{itemize}

Сложность взлома растёт экспоненциально с увеличением длины ключа, поскольку:
\begin{enumerate}
    \item уменьшается количество символов в каждой подгруппе для частотного анализа;
    \item увеличивается число возможных комбинаций;
    \item уменьшается вероятность появления повторяющихся n-грамм.
\end{enumerate}

\section*{3. Какие современные шифры являются наследниками идеи полиалфавитности?}

Современные шифры используют те же принципы смены алфавитов (динамического преобразования), но в более сложной форме:

\begin{itemize}
    \item \textbf{Потоковые шифры (RC4, ChaCha20 и др.)} --- динамический ключевой поток создаёт постоянно меняющуюся таблицу подстановки. Каждый бит/байт шифруется уникальным значением гаммы;

    \item \textbf{Блочные шифры (AES, DES)} --- применяют серию раундов, где каждый раунд фактически формирует свой <<алфавит>> преобразования. S-блоки обеспечивают нелинейную замену;

    \item \textbf{Шифры на основе гаммирования} --- близки идее Виженера, так как используют побитовое сложение (XOR) с ключевым потоком;

    \item \textbf{Режимы работы блочных шифров (CTR, OFB)} --- генерируют псевдослучайный поток для XOR с открытым текстом, что аналогично бесконечному ключу Виженера.
\end{itemize}

Таким образом, идея полиалфавитного шифрования эволюционировала в сложные криптографические примитивы, формирующие меняющиеся преобразования на каждом шаге.
